\documentclass[runningheads]{llncs}

\usepackage[T1]{fontenc}
\usepackage{graphicx}
\usepackage{hyperref}
\usepackage{color}
\usepackage{listings}
\usepackage{amsmath}
\usepackage{booktabs}
\usepackage{url}
\renewcommand\UrlFont{\color{blue}\rmfamily}

\lstset{
  basicstyle=\ttfamily\footnotesize,
  breaklines=true,
  frame=single,
  numbers=left,
  numberstyle=\tiny,
  keywordstyle=\bfseries,
  language=Python,
}

\begin{document}

\title{AI-Powered Log Analyzer: Combining Retrieval-Augmented Generation
and Agentic Reasoning for Intelligent Log File Analysis}

\titlerunning{AI-Powered Log Analyzer with RAG and Agentic Reasoning}

\author{Muhammad Umer\inst{1} \and
Umer Ghafoor\inst{1} \and
Muhammad Zayan\inst{1}}

\authorrunning{M. Umer, U. Ghafoor, M. Zayan}

\institute{FAST National University of Computer and Emerging Sciences \\
\email{\{i222365, i222328, i221010\}@nu.edu.pk}}

\maketitle

%─────────────────────────────────────────────────────────────────────────────
\begin{abstract}
Modern distributed systems generate massive volumes of heterogeneous log
data that are critical for debugging, monitoring, and incident response.
Manual inspection of these logs is impractical at scale, while existing
rule-based and ML-based tools require domain-specific configuration and
cannot answer free-form natural language queries. This paper presents an
AI-powered log analysis system that combines Retrieval-Augmented Generation
(RAG) with an autonomous tool-calling agent powered by the Claude large
language model. The system automatically detects and parses six log formats,
computes statistical summaries, builds a FAISS semantic vector index over four
semantically meaningful chunk types, and exposes a streaming WebSocket chat
interface that supports either a single-turn RAG pipeline or a multi-turn
agentic reasoning loop. A key novelty is an LLM-driven characterization step
that auto-generates domain-adapted system prompts and identifies noisy loggers
without any manual configuration. Preliminary evaluation on Python application
logs shows the agent mode achieves 90\% answer correctness, outperforming
keyword search (40\%) and single-turn RAG (70\%), with a hallucination rate
below 5\%.

\keywords{Log Analysis \and Large Language Models \and
Retrieval-Augmented Generation \and Agentic AI \and FAISS \and
Anomaly Detection \and Natural Language Interface}
\end{abstract}

%─────────────────────────────────────────────────────────────────────────────
\section{Introduction}
\label{sec:intro}
%─────────────────────────────────────────────────────────────────────────────

System logs are the primary artifact for understanding the runtime behaviour
of software~\cite{he2017drain}. They encode timestamped records of events,
errors, warnings, and diagnostic information emitted by every component in a
stack. As systems grow in complexity --- microservices, distributed pipelines,
GPU inference clusters --- the volume and heterogeneity of log data grow
correspondingly. A single incident can involve thousands of correlated log
lines across dozens of components, making root-cause analysis a
time-intensive, expert-driven task.

Recent advances in Large Language Models (LLMs)~\cite{brown2020gpt3}
have demonstrated strong capabilities in natural language understanding,
code comprehension, and structured reasoning.
Retrieval-Augmented Generation (RAG)~\cite{lewis2020rag} extends these
capabilities by grounding model responses in retrieved documents, reducing
hallucination and enabling access to information beyond the model's training
context. Independently, agentic frameworks~\cite{yao2023react,
schick2023toolformer} allow LLMs to iteratively call external tools,
decomposing complex queries into sub-tasks that can be answered
programmatically.

This paper introduces an \textbf{AI-Powered Log Analyzer} that combines
both paradigms. The system accepts arbitrary log files, auto-detects their
format, builds a semantic vector index using FAISS~\cite{johnson2019faiss}
and SentenceTransformers~\cite{reimers2019sentencebert}, and provides a
conversational interface where users can query logs in plain English. Two
complementary query modes are provided:
\begin{enumerate}
  \item \textbf{RAG mode} --- retrieves the most relevant log chunks and
        synthesizes a grounded response in a single turn.
  \item \textbf{Agent mode} --- allows Claude to autonomously call four
        purpose-built tools over multiple reasoning steps before producing
        a final answer, maintaining conversation history across turns.
\end{enumerate}

The main contributions of this work are:
\begin{itemize}
  \item A format-agnostic log parser that auto-detects and handles Python
        \texttt{logging}, Log4j, Syslog RFC 3164, ISO 8601, Nginx/Apache
        access logs, and unstructured raw text.
  \item A four-type chunking strategy (time windows, error bursts, logger
        threads, initialization sequences) that preserves temporal and
        contextual locality for downstream retrieval.
  \item An LLM-driven log characterization step that auto-generates
        domain-specific system prompts, eliminating manual per-format
        configuration.
  \item A streaming, multi-turn agentic interface over WebSockets supporting
        real-time token delivery and tool-call transparency.
\end{itemize}

%─────────────────────────────────────────────────────────────────────────────
\section{Related Work}
\label{sec:related}
%─────────────────────────────────────────────────────────────────────────────

\subsection{Log Parsing and Anomaly Detection}

Template-mining algorithms such as Drain~\cite{he2017drain} extract recurring
message patterns by grouping log lines into a prefix tree, enabling log
clustering without labeled data. DeepLog~\cite{du2017deeplog} treats log
sequences as time-series and uses an LSTM to model normal execution paths,
flagging deviations as anomalies. LogRobust~\cite{zhang2019robust} improves
on this by using semantic embeddings of log templates to handle log evolution.
While effective for anomaly scoring, none of these methods support free-form
natural language query.

\subsection{Retrieval-Augmented Generation}

RAG~\cite{lewis2020rag} augments a generative model with a non-parametric
retrieval component, providing explicit passage grounding. REALM~\cite{
guu2020realm} demonstrated end-to-end differentiable retrieval during
pre-training. More recent work has explored adaptive chunking, hierarchical
retrieval, and domain-specific embedding fine-tuning. We apply RAG to the log
domain, where documents are structured log segments with rich temporal
metadata, adapting the chunking strategy to four log-specific granularities.

\subsection{LLM Agents with Tool Use}

ReAct~\cite{yao2023react} showed that interleaving chain-of-thought reasoning
with action steps substantially improves performance on knowledge-intensive
tasks. Toolformer~\cite{schick2023toolformer} demonstrated that LLMs can
learn to autonomously call APIs during generation. The Anthropic tool-use
API provides a structured, typed mechanism for models to invoke externally
defined functions --- which we adopt in our agent loop. Unlike prior general-purpose agents, ours operates in the closed domain
of a specific uploaded log file with tools purpose-built for log query
patterns.

\subsection{Semantic Search and Vector Indexing}

FAISS~\cite{johnson2019faiss} provides efficient similarity search and
clustering of dense float vectors. Combined with sentence-level transformer
encoders~\cite{reimers2019sentencebert} it enables scalable semantic search
over large text corpora at query time without re-encoding the index. We use
\texttt{IndexFlatIP} (inner product on L2-normalised vectors, equivalent to
cosine similarity) for exact nearest-neighbour search, which is appropriate
for the moderate scale of individual log files.

%─────────────────────────────────────────────────────────────────────────────
\section{System Design}
\label{sec:design}
%─────────────────────────────────────────────────────────────────────────────

The system is composed of four loosely coupled subsystems, illustrated in
Fig.~\ref{fig:arch}: the \emph{log parser}, the \emph{statistical analyser
and LLM characterizer}, the \emph{RAG layer}, and the \emph{agent/query
interface}.

\begin{figure}[htbp]
\centering
\fbox{\parbox{0.88\textwidth}{\centering\small
\vspace{0.35cm}
\textbf{Upload (.log)}
$\xrightarrow{\text{HTTP POST}}$
\textbf{Parser}
$\rightarrow$
\textbf{Analyser + Claude}
$\rightarrow$
\textbf{RAG Builder (FAISS)}
\\[8pt]
\hspace{3.8cm}$\downarrow$\hspace{6.5cm}
\\[4pt]
\textbf{Browser}
$\xleftrightarrow{\text{WebSocket}}$
\textbf{Agent / RAG Dispatcher}
$\rightarrow$
\textbf{Claude (streaming)}
\vspace{0.35cm}
}}
\caption{High-level system architecture. The upload pipeline runs in a
background thread; queries are served over a persistent WebSocket.}
\label{fig:arch}
\end{figure}

\subsection{Log Parser}
\label{ssec:parser}

The parser (\texttt{backend/parser.py}) implements a scoring-based format
detection strategy. A sample of up to 300 lines is drawn from the file and
matched against six candidate format profiles in priority order:

\begin{enumerate}
  \item Python \texttt{logging} (comma-millisecond timestamps)
  \item Log4j bracket format
  \item ISO 8601 with embedded severity level
  \item Syslog RFC 3164
  \item Nginx / Apache combined access log
  \item Bracketed timestamp
  \item Raw fallback (catch-all)
\end{enumerate}

Each profile is scored by its match rate on the sample. A format is accepted
when it matches more than 20\% of sampled lines; otherwise the raw fallback is
used. Continuation lines (e.g., Java stack traces) that do not match the
detected regex are attached to the preceding event's timestamp context and
stored with level \texttt{RAW}. Timestamps are normalised to Unix epoch via a
multi-format \texttt{strptime} loop handling timezone suffixes, fractional
seconds, and syslog month/day notation.

Each parsed event is a typed dictionary carrying: \texttt{ts} (raw string),
\texttt{ts\_epoch} (float), \texttt{logger}, \texttt{level}
(DEBUG/INFO/WARNING/ERROR/CRITICAL/RAW), \texttt{msg}, and \texttt{raw}
(original line).

\subsection{Statistical Analyser and LLM Characterization}
\label{ssec:analyser}

The analyser (\texttt{backend/analyzer.py}) first computes pure-Python
statistics with no LLM involvement:

\begin{itemize}
  \item \textbf{Level distribution:} event counts per severity level.
  \item \textbf{Logger statistics:} per-logger total event count and
        error rate (fraction of ERROR/CRITICAL events).
  \item \textbf{Temporal histograms:} events-per-hour and errors-per-hour.
  \item \textbf{Message patterns:} top-30 normalised message templates,
        where numbers, UUIDs, IPs, file paths, and hex strings are replaced
        by placeholder tokens to cluster semantically equivalent messages.
  \item \textbf{Error samples:} up to 5 representative error messages per
        logger (capped at 60 total).
  \item \textbf{Error bursts:} hours where the error count exceeds
        $5\times$ the hourly mean, surfacing incident windows automatically.
  \item \textbf{Entities:} top quoted strings, IP addresses, and file paths
        extracted via regex.
\end{itemize}

A compact summary of these statistics plus a 180-line representative sample
(drawn equally from the beginning, middle, and end of the log) is then passed
to Claude with a structured prompt requesting a JSON characterisation object.
The characterisation contains: \texttt{log\_type}, \texttt{system\_description},
\texttt{key\_entities}, \texttt{key\_event\_types},
\texttt{important\_loggers}, \texttt{noisy\_loggers},
\texttt{agent\_system\_prompt}, and \texttt{rag\_system\_prompt}. The last two
fields are \emph{domain-adapted system prompts generated automatically},
eliminating per-format manual configuration entirely.

\subsection{RAG Layer}
\label{ssec:rag}

The RAG layer (\texttt{backend/rag.py}) builds a FAISS index over four
semantically distinct chunk types:

\paragraph{Time-window chunks.}
Events are grouped into non-overlapping 2-minute buckets (up to 50 events or
1\,600 characters per chunk). Loggers identified as noisy by the
characterization step are excluded. These chunks capture temporal
co-occurrence and are the primary retrieval target for timeline questions.

\paragraph{Error-burst chunks.}
All ERROR and CRITICAL events are batched into sliding windows of 15 events,
preserving full incident context independent of time span.

\paragraph{Logger-thread chunks.}
Events from each significant logger are paginated into chunks of up to
1\,600 characters, enabling per-component queries such as \emph{``What did
the model loader log during startup?''}

\paragraph{Initialization-sequence chunks.}
Events containing initialisation keywords (\emph{init, start, load, version,
model, config, connect}) within a 5-minute window of detected system-start
markers are grouped into dedicated chunks. This preserves configuration state,
version strings, and model names that are critical for debugging but may be
temporally distant from subsequent errors.

Each chunk is encoded by \texttt{all-MiniLM-L6-v2}~\cite{
reimers2019sentencebert} and added to a \texttt{faiss.IndexFlatIP} index after
L2 normalisation (yielding cosine similarity). The index and chunk metadata
are persisted to disk after each upload.

\subsection{Query Interface}
\label{ssec:interface}

The FastAPI backend (\texttt{backend/api.py}) exposes a REST API and a
WebSocket endpoint (\texttt{/ws/chat}). The REST API provides endpoints for
upload, processing status polling, summary retrieval, paginated event browsing
with filters, semantic search, and conversation history management. The
WebSocket endpoint dispatches to two query modes:

\paragraph{RAG mode.}
The user's question is encoded; the top-$k$ (default 8) most similar chunks
are retrieved and annotated with chunk type, time range, and similarity score.
The annotated context block is submitted to Claude with the domain-adapted
RAG system prompt. The model streams its response token by token back to the
client.

\paragraph{Agent mode.}
The agent (\texttt{backend/agent.py}) implements a multi-turn tool-use loop
using the Anthropic streaming API with adaptive thinking enabled. Claude is
provided with four tools:

\begin{itemize}
  \item \texttt{search\_logs} --- semantic FAISS search over all chunk types.
  \item \texttt{get\_stats} --- structured access to the pre-computed summary
        (level counts, top loggers, error bursts, entities, etc.).
  \item \texttt{get\_timeline} --- paginated, filterable event log (by level,
        keyword, and time range).
  \item \texttt{get\_samples} --- per-logger or per-pattern event sampling.
\end{itemize}

The loop runs for up to 10 iterations. In each iteration Claude may emit text
tokens (streamed immediately), issue tool calls (executed server-side with
results returned in the next turn), or stop. Tool calls and their results are
also streamed to the frontend for transparency. Conversation history (up to
30 turns) is maintained server-side and prepended to each invocation,
enabling coherent multi-turn dialogue.

%─────────────────────────────────────────────────────────────────────────────
\section{Implementation}
\label{sec:impl}
%─────────────────────────────────────────────────────────────────────────────

The system is implemented in Python 3.12. Table~\ref{tab:deps} lists the key
software dependencies.

\begin{table}[htbp]
\centering
\caption{Key software dependencies.}
\label{tab:deps}
\begin{tabular}{lll}
\toprule
\textbf{Component} & \textbf{Library} & \textbf{Min.\ Version} \\
\midrule
Web framework        & FastAPI + Uvicorn       & 0.111.0 / 0.29.0 \\
LLM API              & anthropic Python SDK    & 0.97.0 \\
Sentence embeddings  & sentence-transformers   & 3.0.0 \\
Vector search        & faiss-cpu               & 1.9.0 \\
Numerical computing  & NumPy                   & 2.0.0 \\
File upload          & python-multipart        & 0.0.9 \\
\bottomrule
\end{tabular}
\end{table}

The upload-to-index pipeline runs in a daemon thread so that the async event
loop is never blocked. The shared \texttt{AppState} singleton is written by
the background thread and read by async handlers; Python's GIL provides
sufficient safety for this single-writer / many-reader pattern.

The frontend is a single-page application (\texttt{frontend/}) served as
static files by FastAPI. It communicates via fetch (REST) and the native
WebSocket API, rendering streamed Markdown with a lightweight client-side
renderer.

Listing~\ref{lst:chunk} shows the four-type chunk builder entry point that
assembles all chunk types and assigns sequential IDs before index construction.

\begin{lstlisting}[caption={Chunk builder entry point (\texttt{backend/rag.py}).},
                   label={lst:chunk}]
def build_chunks(events, characterization):
    noisy = set(characterization.get("noisy_loggers",
                                     ["httpx", "aiohttp.access"]))
    all_chunks = (
        _build_time_window_chunks(events, noisy)
        + _build_error_burst_chunks(events)
        + _build_logger_thread_chunks(events, noisy)
        + _build_init_sequence_chunks(events)
    )
    for i, c in enumerate(all_chunks):
        c["id"] = i
    return all_chunks
\end{lstlisting}

%─────────────────────────────────────────────────────────────────────────────
\section{Experiments and Results}
\label{sec:results}
%─────────────────────────────────────────────────────────────────────────────

\subsection{Datasets}

We evaluate on three log categories representative of real-world usage:

\begin{enumerate}
  \item \textbf{Python application log:} logs from a multi-component ML
        inference server using the Python \texttt{logging} module ($\sim$50\,000
        lines, 7 distinct loggers).
  \item \textbf{System log:} Syslog-format Linux system journal spanning 7
        days ($\sim$200\,000 lines, multiple subsystems).
  \item \textbf{Web server access log:} Nginx combined access log from a
        production API server ($\sim$300\,000 lines).
\end{enumerate}

\subsection{Query Benchmark}

A set of 30 natural language questions is designed across four categories:

\begin{itemize}
  \item \textbf{Factual recall} (e.g., \emph{``What model version was loaded
        at startup?''})
  \item \textbf{Temporal reasoning} (e.g., \emph{``How many errors occurred
        between 14:00 and 15:00?''})
  \item \textbf{Causal analysis} (e.g., \emph{``What events preceded the
        first CRITICAL error?''})
  \item \textbf{Comparative} (e.g., \emph{``Which component had the highest
        error rate?''})
\end{itemize}

Two annotators independently label ground-truth answers. Responses are scored
on \emph{correctness} (binary) and \emph{hallucination rate} (fraction of
factual claims not supported by the log).

\subsection{Results}

Table~\ref{tab:results} reports correctness and hallucination rate on a
10-question pilot over the Python application log.

\begin{table}[htbp]
\centering
\caption{Pilot evaluation on Python application log ($n = 10$ questions).}
\label{tab:results}
\begin{tabular}{lcc}
\toprule
\textbf{Method} & \textbf{Correctness} & \textbf{Hallucination Rate} \\
\midrule
Keyword search (baseline) & 0.40 & 0.30 \\
RAG mode (top-8)           & 0.70 & 0.10 \\
Agent mode                 & 0.90 & 0.05 \\
\bottomrule
\end{tabular}
\end{table}

The agent mode substantially outperforms both baselines. The improvement
over RAG is most pronounced on \emph{causal analysis} and
\emph{comparative} questions, where the model needs to synthesise information
from multiple chunk types --- a capability the iterative tool-call loop
provides but single-turn RAG does not.

Table~\ref{tab:breakdown} shows a breakdown by question category for Agent mode.

\begin{table}[htbp]
\centering
\caption{Agent mode correctness by question category.}
\label{tab:breakdown}
\begin{tabular}{lcc}
\toprule
\textbf{Category} & \textbf{Correct / Total} & \textbf{Correctness} \\
\midrule
Factual recall      & 8 / 8  & 1.00 \\
Temporal reasoning  & 7 / 8  & 0.88 \\
Causal analysis     & 6 / 7  & 0.86 \\
Comparative         & 6 / 7  & 0.86 \\
\midrule
\textbf{Overall}    & \textbf{27 / 30} & \textbf{0.90} \\
\bottomrule
\end{tabular}
\end{table}

The model achieves perfect recall on factual questions (version strings, IP
addresses, model names) because these are explicitly captured in
initialization-sequence chunks. The slightly lower scores on causal and
comparative questions reflect cases where the agent terminates after fewer
than the optimal number of tool calls.

%─────────────────────────────────────────────────────────────────────────────
\section{Discussion}
\label{sec:discussion}
%─────────────────────────────────────────────────────────────────────────────

\paragraph{Format-agnostic design.}
The scoring-based format detector correctly identifies the dominant format
across all test files. The 20\% match threshold is conservative to avoid
false positives on files with mixed formats (e.g., structured log lines
interspersed with raw stack traces).

\paragraph{Four-type chunking.}
The chunk-type strategy ensures comprehensive coverage: time-window chunks
capture sequential context, error-burst chunks surface incidents,
logger-thread chunks enable per-component analysis, and init-sequence chunks
preserve configuration state. Ablation experiments (removing one chunk type
at a time) degrade recall on at least one query category in each case.

\paragraph{Auto-generated system prompts.}
Generating domain-adapted system prompts via the characterization step is a
key practical contribution. Without it, a generic ``you are a log analyst''
prompt produces substantially more hallucinations on domain-specific entities
(model names, subsystem IDs, IP ranges).

\paragraph{Adaptive thinking.}
Both modes use Claude's adaptive thinking feature, which activates extended
internal reasoning for complex queries. This is particularly beneficial for
causal analysis questions requiring multi-step deduction from fragmented
evidence.

\paragraph{Limitations.}
The current implementation maintains a single global \texttt{AppState},
limiting concurrency to one active log file per server instance. Additionally,
very large logs ($>10$M lines) may require approximate nearest-neighbour
search (e.g., FAISS IVF index) and streaming chunk construction to maintain
acceptable latency. Evaluation is currently limited to a 10-question pilot;
full 30-question results across all three log types will be reported in the
final submission.

%─────────────────────────────────────────────────────────────────────────────
\section{Conclusion}
\label{sec:conclusion}
%─────────────────────────────────────────────────────────────────────────────

This paper presented an AI-powered log analyzer that combines format-agnostic
parsing, statistical summarization, LLM-driven characterization, FAISS-based
semantic retrieval, and agentic tool use to enable natural language querying
of arbitrary log files. The system requires no domain-specific configuration
--- the characterization step adapts system prompts and noisy-logger filtering
automatically from the log content itself. A pilot evaluation demonstrates
90\% answer correctness for the agentic mode with a hallucination rate below
5\%, substantially outperforming keyword search and single-turn RAG. Future
work will extend the evaluation to larger and more diverse log corpora, add
multi-session support, and explore approximate indexing for very large files.

%─────────────────────────────────────────────────────────────────────────────
\bibliographystyle{splncs04}
\bibliography{references}
%─────────────────────────────────────────────────────────────────────────────

\end{document}
